\documentclass[../../main.tex]{subfiles}

\begin{document}

\section{有限测度集上的有界可测函数的 Lebesgue 积分}\label{section:ch04-有界可测积分}

我们前一节讨论的 Dirichlet 函数显现了 Riemann 积分的一个主要缺陷: 在闭有界区间上的一致有界的 Riemann 可积函数序列可能逐点收敛到不是 Riemann 可积的函数. 我们将看到 Lebesgue 积分没有这个缺陷.

今后除非直接提到, 否则我们仅考虑 Lebesgue 积分, 因而我们用纯积分号表示 Lebesgue 积分. 即将介绍的定理 3 告诉我们任何在 $[a,b]$ 上是 Riemann 可积的有界函数在 $[a,b]$ 上也是 Lebesgue 可积的, 且这两个积分相等.

回忆定义在集合 $E$ 上的一个可测的实值函数 $\psi$ 称为是简单的, 若它仅取有限个实值. 若 $\psi$ 在 $E$ 上取不同值 $a_{1},\cdots,a_{n}$, 则根据 $\psi$ 的可测性, 它的水平集 $\psi^{-1}(a_{i})$ 是可测的且我们有 $\psi$ 在 $E$ 上的典范表示:
\begin{align}
\text{在 }E\text{ 上},\quad\psi=\sum_{i=1}^{n}a_{i}\cdot\chi_{E_{i}},\quad\text{其中每个 }E_{i}=\psi^{-1}(a_{i})=\{x\in E\mid\psi(x)=a_{i}\}.\label{align:ch04-简单典范}
\end{align}
典范表示的特征是 $E_{i}$ 是不交的且 $a_{i}$ 是不同的.

\begin{definition}[简单函数的积分]\label{definition:ch04-简单积分}
对于定义在有限测度集 $E$ 上的简单函数 $\psi$, 我们定义 $\psi$ 在 $E$ 上的积分为
\begin{align*}
\int_{E}\psi=\sum_{i=1}^{n}a_{i}\cdot m(E_{i}),
\end{align*}
其中 $\psi$ 具有由\eqref{align:ch04-简单典范}给出的典范表示.
\end{definition}

\begin{lemma}\label{lemma:ch04-简单积分良定}
令 $\{E_{i}\}_{i=1}^{n}$ 为有限测度集 $E$ 的可测子集的有限不交族. 对于 $1\leqslant i\leqslant n$, 令 $a_{i}$ 为实数,
\begin{align*}
\text{若在 }E\text{ 上 }\varphi=\sum_{i=1}^{n}a_{i}\cdot\chi_{E_{i}},\quad\text{则}\int_{E}\varphi=\sum_{i=1}^{n}a_{i}\cdot m(E_{i}).
\end{align*}
\end{lemma}

\begin{proof}
族 $\{E_{i}\}_{i=1}^{n}$ 是不交的但上面的式子可以不是典范表示, 这是因为这些 $a_{i}$ 可以相同. 我们必须考虑可能的重复. 令 $\{\lambda_{1},\cdots,\lambda_{m}\}$ 为 $\varphi$ 取的不同值. 对于 $1\leqslant j\leqslant m$, 设 $A_{j}=\{x\in E\mid\varphi(x)=\lambda_{j}\}$. 根据通过典范表示给出的积分的定义,
\begin{align*}
\int_{E}\varphi=\sum_{j=1}^{m}\lambda_{j}\cdot m(A_{j}).
\end{align*}
对于 $1\leqslant j\leqslant m$, 令 $I_{j}$ 为 $\{1,\cdots,n\}$ 中那些使得 $a_{i}=\lambda_{j}$ 的指标 $i$ 构成的集合. 则 $\{1,\cdots,n\}=\bigcup_{j=1}^{m}I_{j}$, 且这个并是不交的. 此外, 根据测度的有限可加性,
\begin{align*}
\text{对所有 }1\leqslant j\leqslant m,\quad m(A_{j})=\sum_{i\in I_{j}}m(E_{i}).
\end{align*}
因此
\begin{align*}
\sum_{i=1}^{n}a_{i}\cdot m(E_{i})=\sum_{j=1}^{m}\left[\sum_{i\in I_{j}}a_{i}\cdot m(E_{i})\right]=\sum_{j=1}^{m}\lambda_{j}\left[\sum_{i\in I_{j}}m(E_{i})\right]=\sum_{j=1}^{m}\lambda_{j}\cdot m(A_{j})=\int_{E}\varphi.
\end{align*}
\end{proof}

我们的目标之一是对一般的 Lebesgue 积分建立线性与单调性质. 以下是我们在这个方向的第一个结果.

\begin{proposition}[积分的线性与单调性]\label{proposition:ch04-简单积分线性}
令 $\varphi$ 和 $\psi$ 为定义在有限测度集 $E$ 上的简单函数. 则对任何 $\alpha$ 和 $\beta$,
\begin{align*}
\int_{E}(\alpha\varphi+\beta\psi)=\alpha\int_{E}\varphi+\beta\int_{E}\psi.
\end{align*}
此外,
\begin{align*}
\text{若在 }E\text{ 上 }\varphi\leqslant\psi,\quad\text{则}\int_{E}\varphi\leqslant\int_{E}\psi.
\end{align*}
\end{proposition}

\begin{proof}
由于 $\varphi$ 和 $\psi$ 在 $E$ 上都仅取有限个值, 我们可以选取 $E$ 的可测子集的有限不交族 $\{E_{i}\}_{i=1}^{n}$, 该族的并是 $E$, 使得 $\varphi$ 和 $\psi$ 在每个 $E_{i}$ 上是常数. 对每个 $i$ ($1\leqslant i\leqslant n$), 令 $a_{i}$ 和 $b_{i}$ 分别为 $\varphi$ 和 $\psi$ 在 $E_{i}$ 上所取的值. 根据前一个引理,
\begin{align*}
\int_{E}\varphi=\sum_{i=1}^{n}a_{i}\cdot m(E_{i})\quad\text{且}\quad\int_{E}\psi=\sum_{i=1}^{n}b_{i}\cdot m(E_{i}).
\end{align*}
然而, 简单函数 $\alpha\varphi+\beta\psi$ 在 $E_{i}$ 上取常数值 $\alpha a_{i}+\beta b_{i}$. 因此, 又一次由前一个引理,
\begin{align*}
\int_{E}(\alpha\varphi+\beta\psi)=\sum_{i=1}^{n}(\alpha a_{i}+\beta b_{i})\cdot m(E_{i})=\alpha\sum_{i=1}^{n}a_{i}\cdot m(E_{i})+\beta\sum_{i=1}^{n}b_{i}\cdot m(E_{i})=\alpha\int_{E}\varphi+\beta\int_{E}\psi.
\end{align*}
为证明单调性, 假设在 $E$ 上 $\varphi\leqslant\psi$. 在 $E$ 上定义 $\eta=\psi-\varphi$. 由线性性质,
\begin{align*}
\int_{E}\psi-\int_{E}\varphi=\int_{E}(\psi-\varphi)=\int_{E}\eta\geqslant 0,
\end{align*}
这是因为非负简单函数 $\eta$ 有一个非负积分.
\end{proof}

有限测度集上的简单函数的积分的线性表明, 在引理 1 的陈述中族 $\{E_{i}\}_{i=1}^{n}$ 是不交的限制是不必要的.

阶梯函数仅取有限个值而每个区间是可测的. 因此阶梯函数是简单的. 由于单点集的测度是零而区间的测度是它的长度, 我们从定义在有限测度集上的简单函数的 Lebesgue 积分的线性推出闭有界区间上的阶梯函数的 Riemann 积分与 Lebesgue 积分一致.

令 $f$ 为定义在有限测度集 $E$ 上的有界实值函数. 类似于 Riemann 积分, 我们分别定义 $f$ 在 $E$ 上的下 Lebesgue 积分与上 Lebesgue 积分为
\begin{align*}
\sup\left\{\int_{E}\varphi\mid\varphi\text{ 是简单的且在 }E\text{ 上 }\varphi\leqslant f\right\}
\end{align*}
与
\begin{align*}
\inf\left\{\int_{E}\psi\mid\psi\text{ 是简单的且在 }E\text{ 上 }f\leqslant\psi\right\}.
\end{align*}
由于假设 $f$ 有界, 根据简单函数的积分的单调性质, 下积分和上积分是有限的且上积分总是至少与下积分一样大.

\begin{definition}[Lebesgue 可积]\label{definition:ch04-Lebesgue可积}
在有限测度集 $E$ 上的有界函数 $f$ 称为在 $E$ 上是 Lebesgue 可积的, 若它在 $E$ 上的上 Lebesgue 积分与下 Lebesgue 积分相等. 上积分与下积分的公共值称为 $f$ 在 $E$ 上的 Lebesgue 积分, 或简单地称为 $f$ 在 $E$ 上的积分, 记为 $\int_{E}f$.
\end{definition}

\begin{theorem}\label{theorem:ch04-Riemann蕴含Lebesgue}
令 $f$ 为定义在闭有界区间 $[a,b]$ 上的有界函数. 若 $f$ 在 $[a,b]$ 上是 Riemann 可积的, 则它在 $[a,b]$ 上是 Lebesgue 可积的且这两个积分是相等的.
\end{theorem}

\begin{proof}
$f$ 是 Riemann 可积的断言意味着, 设 $I=[a,b]$,
\begin{align*}
\sup\left\{(R)\int_{I}\varphi\mid\varphi\text{ 是阶梯函数},\varphi\leqslant f\right\}=\inf\left\{(R)\int_{I}\psi\mid\psi\text{ 是阶梯函数},f\leqslant\psi\right\}.
\end{align*}
为证明 $f$ 是 Lebesgue 可积的, 我们必须证明:
\begin{align*}
\sup\left\{\int_{I}\varphi\mid\varphi\text{ 是简单的},\varphi\leqslant f\right\}=\inf\left\{\int_{I}\psi\mid\psi\text{ 是简单的},f\leqslant\psi\right\}.
\end{align*}
然而, 每个阶梯函数是简单函数, 我们已观察到, 对于阶梯函数, Riemann 积分和 Lebesgue 积分是相同的. 因此第一个等式蕴涵第二个等式且蕴涵 Riemann 积分和 Lebesgue 积分的相等性.
\end{proof}

我们现在完全证使用符号 $\int_{E}f$ 表示有限测度集上的 Lebesgue 可积的有界函数的积分的合理性, 无须任何关于 $(R)$ 的前缀. 在区间 $E=[a,b]$ 的情形, 我们有时用熟悉的记号 $\int_{a}^{b}f$ 表示 $\int_{[a,b]}f$, 而有时用经典的 Leibniz 记号 $\int_{a}^{b}f(x)\,dx$ 是有用的.

\begin{example}\label{example:ch04-Dirichlet-Lebesgue}
$[0,1]$ 中的有理数集 $E$ 是测度为零的可测集. Dirichlet 函数 $f$ 是 $E$ 的特征函数 $\chi_{E}$ 在 $[0,1]$ 上的限制. 因此 $f$ 在 $[0,1]$ 上可积且
\begin{align*}
\int_{[0,1]}f=\int_{[0,1]}1\cdot\chi_{E}=1\cdot m(E)=0.
\end{align*}
我们已证明 $f$ 在 $[0,1]$ 上不是 Riemann 可积的.
\end{example}

\begin{theorem}\label{theorem:ch04-有界可测可积}
令 $f$ 为有限测度集 $E$ 上的有界可测函数. 则 $f$ 在 $E$ 上是可积的.
\end{theorem}

\begin{proof}
令 $n$ 为自然数. 根据简单逼近引理, 其中 $\varepsilon=1/n$, 存在定义在 $E$ 上的两个简单函数 $\varphi_{n}$ 与 $\psi_{n}$, 使得
\begin{align*}
\text{在 }E\text{ 上},\quad\varphi_{n}\leqslant f\leqslant\psi_{n}
\end{align*}
且
\begin{align*}
\text{在 }E\text{ 上},\quad 0\leqslant\psi_{n}-\varphi_{n}\leqslant 1/n.
\end{align*}
根据简单函数的积分的单调性与线性,
\begin{align*}
0\leqslant\int_{E}\psi_{n}-\int_{E}\varphi_{n}=\int_{E}[\psi_{n}-\varphi_{n}]\leqslant 1/n\cdot m(E).
\end{align*}
然而,
\begin{align*}
0\leqslant\inf\left\{\int_{E}\psi\mid\psi\text{ 是简单的},\psi\geqslant f\right\}-\sup\left\{\int_{E}\varphi\mid\varphi\text{ 是简单的},\varphi\leqslant f\right\}\leqslant\int_{E}\psi_{n}-\int_{E}\varphi_{n}\leqslant 1/n\cdot m(E).
\end{align*}
该不等式对每个自然数 $n$ 成立, 而 $m(E)$ 是有限的. 因此上 Lebesgue 积分与下 Lebesgue 积分相等, 因此函数 $f$ 在 $E$ 上可积.
\end{proof}

事实证明前一个定理的逆命题成立, 有限测度集上的有界函数是 Lebesgue 可积的当且仅当它是可测的. 我们稍后证明（见第 5 章的\cref{theorem:ch05-Riemann刻画}）这个定理. 特别地, 这表明不是每个定义在有限测度集上的有界函数都是 Lebesgue 可积的. 事实上, 对任何有限正测度的可测集 $E$, 其每个不可测子集的特征函数在 $E$ 上的限制不是 Lebesgue 可积的.

\begin{theorem}[积分的线性与单调性]\label{theorem:ch04-有界积分线性}
令 $f$ 和 $g$ 为有限测度集 $E$ 上的有界可测函数. 则对任何 $\alpha$ 和 $\beta$,
\begin{align}
\int_{E}(\alpha f+\beta g)=\alpha\int_{E}f+\beta\int_{E}g.\label{align:ch04-线性}
\end{align}
此外,
\begin{align}
\text{若在 }E\text{ 上 }f\leqslant g,\text{ 则}\int_{E}f\leqslant\int_{E}g.\label{align:ch04-单调性}
\end{align}
\end{theorem}

\begin{proof}
可测有界函数的线性组合也是可测与有界的. 因此, 根据\cref{theorem:ch04-有界可测可积}, $\alpha f+\beta g$ 在 $E$ 上是可积的. 我们首先对 $\beta=0$ 证明线性. 若 $\psi$ 是简单函数, $\alpha\psi$ 也是, 反过来（若 $\alpha\neq 0$）也如此. 我们对简单函数证明积分的线性. 令 $\alpha>0$. 由于 Lebesgue 积分等于上 Lebesgue 积分,
\begin{align*}
\int_{E}\alpha f=\inf_{\psi\geqslant\alpha f}\int_{E}\psi=\alpha\inf_{[\psi/\alpha]\geqslant f}\int_{E}[\psi/\alpha]=\alpha\int_{E}f.
\end{align*}
对于 $\alpha<0$, 由于 Lebesgue 积分既等于上 Lebesgue 积分又等于下 Lebesgue 积分,
\begin{align*}
\int_{E}\alpha f=\inf_{\psi\geqslant\alpha f}\int_{E}\psi=\alpha\sup_{[\psi/\alpha]\leqslant f}\int_{E}[\varphi/\alpha]=\alpha\int_{E}f.
\end{align*}
剩下来要在 $\alpha=\beta=1$ 的情形下证明线性. 令 $\psi_{1}$ 和 $\psi_{2}$ 为满足在 $E$ 上 $f\leqslant\psi_{1}$ 且 $g\leqslant\psi_{2}$ 的简单函数, 则 $\psi_{1}+\psi_{2}$ 是一个简单函数且在 $E$ 上, $f+g\leqslant\psi_{1}+\psi_{2}$. 因此, 由于 $\int_{E}(f+g)$ 等于 $f+g$ 在 $E$ 上的上 Lebesgue 积分, 根据简单函数的积分的线性,
\begin{align*}
\int_{E}(f+g)\leqslant\int_{E}(\psi_{1}+\psi_{2})=\int_{E}\psi_{1}+\int_{E}\psi_{2}.
\end{align*}
当 $\psi_{1}$ 和 $\psi_{2}$ 在满足 $f\leqslant\psi_{1}$ 和 $g\leqslant\psi_{2}$ 的简单函数中变动时, 右边积分和的最大下界等于 $\int_{E}f+\int_{E}g$. 这些不等式告诉我们 $\int_{E}(f+g)$ 是这些和的一个下界. 因此,
\begin{align*}
\int_{E}(f+g)\leqslant\int_{E}f+\int_{E}g.
\end{align*}
剩下来要证明反向的不等式. 令 $\varphi_{1}$ 和 $\varphi_{2}$ 为满足在 $E$ 上 $\varphi_{1}\leqslant f$ 且 $\varphi_{2}\leqslant g$ 的简单函数. 则在 $E$ 上 $\varphi_{1}+\varphi_{2}\leqslant f+g$ 且 $\varphi_{1}+\varphi_{2}$ 是简单的. 因此, 由于 $\int_{E}(f+g)$ 等于 $f+g$ 在 $E$ 上的下 Lebesgue 积分, 根据简单函数的积分的线性,
\begin{align*}
\int_{E}(f+g)\geqslant\int_{E}(\varphi_{1}+\varphi_{2})=\int_{E}\varphi_{1}+\int_{E}\varphi_{2}.
\end{align*}
当 $\varphi_{1}$ 和 $\varphi_{2}$ 在满足 $\varphi_{1}\leqslant f$ 和 $\varphi_{2}\leqslant g$ 的简单函数中变动时, 右边积分和的最小上界等于 $\int_{E}f+\int_{E}g$. 这些不等式告诉我们 $\int_{E}(f+g)$ 是这些相同和的一个上界. 因此,
\begin{align*}
\int_{E}(f+g)\geqslant\int_{E}f+\int_{E}g.
\end{align*}
这完成了积分的线性的证明.

为了证明单调性, 假设在 $E$ 上 $f\leqslant g$. 在 $E$ 上定义 $h=g-f$. 根据线性,
\begin{align*}
\int_{E}g-\int_{E}f=\int_{E}(g-f)=\int_{E}h.
\end{align*}
函数 $h$ 是非负的且因此在 $E$ 上 $\psi\leqslant h$, 其中在 $E$ 上 $\psi\equiv 0$. 由于 $h$ 的积分等于它的下积分, $\int_{E}h\geqslant\int_{E}\psi=0$. 因此 $\int_{E}f\leqslant\int_{E}g$.
\end{proof}

\begin{corollary}\label{corollary:ch04-积分可加}
令 $f$ 为有限测度集 $E$ 上的有界可测函数. 假定 $A$ 和 $B$ 是 $E$ 的不交可测子集. 则
\begin{align}
\int_{A\cup B}f=\int_{A}f+\int_{B}f.\label{align:ch04-可加}
\end{align}
\end{corollary}

\begin{proof}
$f\cdot\chi_{A}$ 和 $f\cdot\chi_{B}$ 都是 $E$ 上的有界可测函数. 由于 $A$ 和 $B$ 不交,
\begin{align*}
f\cdot\chi_{A\cup B}=f\cdot\chi_{A}+f\cdot\chi_{B}.
\end{align*}
此外, 对于 $E$ 的任何可测子集 $E_{1}$（见习题 10）,
\begin{align*}
\int_{E_{1}}f=\int_{E}f\cdot\chi_{E_{1}}.
\end{align*}
因此, 根据积分的线性,
\begin{align*}
\int_{A\cup B}f=\int_{E}f\cdot\chi_{A\cup B}=\int_{E}f\cdot\chi_{A}+\int_{E}f\cdot\chi_{B}=\int_{A}f+\int_{B}f.
\end{align*}
\end{proof}

\begin{corollary}\label{corollary:ch04-积分绝对值}
令 $f$ 为有限测度集 $E$ 上的有界可测函数. 则
\begin{align}
\left|\int_{E}f\right|\leqslant\int_{E}|f|.\label{align:ch04-绝对值}
\end{align}
\end{corollary}

\begin{proof}
函数 $|f|$ 是可测与有界的. 现在
\begin{align*}
\text{在 }E\text{ 上},\quad-|f|\leqslant f\leqslant|f|.
\end{align*}
根据积分的线性与单调性,
\begin{align*}
-\int_{E}|f|\leqslant\int_{E}f\leqslant\int_{E}|f|,
\end{align*}
即\eqref{align:ch04-绝对值}成立.
\end{proof}

\begin{proposition}\label{proposition:ch04-一致收敛积分}
令 $\{f_{n}\}$ 为有限测度集 $E$ 上的有界可测函数序列.
\begin{align*}
\text{若在 }E\text{ 上一致地 }\{f_{n}\}\to f,\quad\text{则}\lim_{n\to\infty}\int_{E}f_{n}=\int_{E}f.
\end{align*}
\end{proposition}

\begin{proof}
由于收敛是一致的且每个 $f_{n}$ 是有界的, 极限函数 $f$ 是有界的. 函数 $f$ 是可测的, 这是由于它是可测函数序列的逐点极限. 令 $\varepsilon>0$. 选取指标 $N$ 使得
\begin{align}
\text{对所有 }n\geqslant N,\quad\text{在 }E\text{ 上 }|f-f_{n}|<\varepsilon/m(E).\label{align:ch04-一致收敛估计}
\end{align}
根据积分的线性与单调性以及前一个推论, 对每个 $n\geqslant N$,
\begin{align*}
\left|\int_{E}f-\int_{E}f_{n}\right|=\left|\int_{E}[f-f_{n}]\right|\leqslant\int_{E}|f-f_{n}|\leqslant[\varepsilon/m(E)]\cdot m(E)=\varepsilon.
\end{align*}
\end{proof}

\end{document}